\documentclass[12pt]{article}
\usepackage[utf8]{inputenc}
\usepackage[T1]{fontenc}
\usepackage{hyperref}
\title{Kele Language Resources for K-12 Curriculum Development}
\author{Compiled by Claude Sonnet 5.5 (claude-sonnet-5-5) for Andr\'as Kornai}
\date{}
\begin{document}
\maketitle

\section{Introduction}
Kele (also Gele'; ISO 639-3 sbc; Glottolog kele1258) is an Austronesian (Oceanic, Admiralty Islands, East Manus) language of inland Manus Island, Papua New Guinea. Sources found give about 1,500 speakers (2020), which is higher than the 600 in the request. Kele is not to be confused with Lele (lle), a separate East Manus language. Very little educational material was found in web searches; no primer, school programme or Kele scripture was located.

\section{Primary Educational Resources}
No Kele primers, textbooks or school programmes were found in search results. The Wikipedia article and the Joshua Project page give only general background on the people and language.
\begin{itemize}
\item Kele language (New Guinea), Wikipedia (English overview): \url{https://en.wikipedia.org/wiki/Kele_language_(New_Guinea)}
\item Gele in Papua New Guinea, Joshua Project: \url{https://joshuaproject.net/people_groups/11858}
\end{itemize}

\section{Linguistic and Reference Materials}
\begin{itemize}
\item University of Queensland manuscripts item F3610: comparative Admiralty Islands wordlists (W.E. Smythe, compiled by John Z'graggen, 2000), listed as accompanied by ``Admiralty Island Linguistics'' Part 4, Gele' Grammar. Availability beyond the catalogue entry was not verified: \url{https://manuscripts.library.uq.edu.au/index.php/f3610}
\item Kele page of the Trans-New-Guinea / Oceanic lexical database (TLOPO): \url{https://tlopo.clld.org/languages/kele}
\item Glottolog entry for Kele: \url{https://glottolog.org/resource/languoid/id/kele1258.bigmap.html}
\item Kele numerals (Max Planck Institute, numeral systems of the world): \url{https://lingweb.eva.mpg.de/channumerals/Kele.htm}
\item Manus languages, Wikipedia (context on the dialect continuum): \url{https://en.wikipedia.org/wiki/Manus_languages}
\end{itemize}

\section{Religious and Community Resources}
No Kele scripture or community organisation was found. The Joshua Project page above is the only community-level entry located.

\section{Digital and Audio Resources}
No Kele audio, video, corpus or software resources were found in search results.

\section{Curriculum Development Context}
Kele is spoken alongside neighbouring Ere and Kurti, and many speakers are multilingual. Documentation is thin (a grammar volume and wordlists are catalogued but not confirmed to be accessible), so any Kele-medium curriculum would first need community orthography decisions, a lexicon and fieldwork to produce reading materials. Fieldwork with the community, and a search of archives such as PARADISEC and SIL, are the recommended next steps; these were not searched here.

\end{document}